% Week 5 progress deck for the QSVM track. Standalone.
%
% Build:  latexmk -pdf qsvm_week5_slides.tex
%
% Figures come from figures_slides/, NOT figures/ — same data, larger type,
% sized for a 16:9 slide. Generate with
%   uv run python scripts/make_report_figures.py --target slides

\documentclass[aspectratio=169,11pt]{beamer}

\usepackage[T1]{fontenc}
\usepackage[utf8]{inputenc}
\usepackage{lmodern}
\usepackage{amsmath,amssymb}
\usepackage{booktabs}
\usepackage{graphicx}
\usepackage{siunitx}
\usepackage{microtype}

% Without detect-weight, siunitx strips \bfseries inside S columns and the
% "best model in bold" markup renders as ordinary text.
\sisetup{detect-weight=true,detect-inline-weight=math}

% Plain frames, no navigation furniture. The content is dense enough that
% sidebars and headline bars would only steal width.
\usetheme{default}
\usecolortheme{default}
\setbeamertemplate{navigation symbols}{}
\setbeamertemplate{itemize items}[circle]
\setbeamertemplate{footline}{%
  \hfill\usebeamerfont{footline}\color{gray!80}\small\insertframenumber\hspace{2em}\vspace{1.5ex}}
\setbeamercolor{frametitle}{fg=black}
\setbeamerfont{frametitle}{size=\large,series=\bfseries}
\setbeamercolor{structure}{fg=blue!55!black}
\setbeamerfont{block title}{size=\normalsize}

\newcommand{\nc}{\ensuremath{n_{\mathrm{c}}}}
\newcommand{\FIG}[1]{figures_slides/#1}
\newcommand{\take}[1]{%
  \vspace{0.4em}\footnotesize\textbf{\color{blue!55!black}Takeaway:} #1}

\title{QSVM Track --- Week 5}
\subtitle{Four-model baselines, kernel diagnostics, and the first significance tests}
\date{2 August 2026}

\begin{document}

\begin{frame}[plain]
  \titlepage
\end{frame}

% ---------------------------------------------------------------- overview
\begin{frame}{Week 5 in one slide}
\small
Three results change conclusions this track had already drawn. They do not
merely add to them.
\vspace{0.6em}

\begin{tabular}{@{}p{0.40\textwidth}p{0.13\textwidth}p{0.42\textwidth}@{}}
\toprule
\textbf{Position entering the week} & \textbf{Now} & \textbf{Basis} \\
\midrule
Classical comparison is SVM only & Resolved & Four models, both tasks \\
CIC 15-class collapse unexplained & \textbf{Attributed} & Kernel concentration vs.\ an RBF control \\
No significance testing & Resolved & 84 paired tests, all $p<0.05$ \\
Subsample fidelity assumed & Verified & KS rejections far below chance \\
SOREL blocked & Closed & Resume conditions recorded \\
\bottomrule
\end{tabular}

\take{the classical--quantum gap is real, larger than we had reported, and we
now know \emph{why} it appears on one dataset and not the other.}
\end{frame}

% ---------------------------------------------------------------- data
\begin{frame}{Data and protocol}
\small
\begin{columns}[T,onlytextwidth]
\begin{column}{0.48\textwidth}
\textbf{CIC-MalMem-2022} --- memory dumps, 55 forensic features
\begin{itemize}\small
  \item binary: \num{58596} rows, balanced
  \item 15-class: \num{29298} malware rows, $1.7\times$ imbalanced
\end{itemize}
\vspace{0.5em}
\textbf{EMBER 2018} --- static PE, 2381 features
\begin{itemize}\small
  \item binary: \num{18014} rows, $88.9\%$ malware
  \item 15-class: \num{14325} rows, 955 per family (balanced)
\end{itemize}
\end{column}
\begin{column}{0.48\textwidth}
\textbf{Fair-comparison protocol}
\begin{itemize}\small
  \item quantum kernel costs $O(n^2)$ circuit evaluations $\Rightarrow$
        stratified 1000-row subsample
  \item that subsample and its 5 outer folds are \emph{persisted and replayed}
        by the classical runner
  \item one shared pipeline: variance filter $\to$ correlation filter $\to$
        standardise $\to$ PCA to \nc{}
  \item \nc{} fixes classical feature count \emph{and} qubit budget at once
  \item fit on training-fold rows only
\end{itemize}
\end{column}
\end{columns}

\vspace{0.4em}
\footnotesize
\textbf{Metric: macro-F1, not accuracy.} On EMBER binary at \nc{}~$=1$ the QSVM
reaches $0.890$ accuracy in 3 of 5 folds while benign-class F1 is exactly
$0.000$ --- $88.9\%$ of rows are malware and the model labels everything malware.
\end{frame}

% ---------------------------------------------------------------- baseline
\begin{frame}{Result 1 --- the classical baseline was too weak}
\small
Every EMBER margin on record compared the QSVM against \emph{SVM alone}.
Adding random forest, XGBoost and LightGBM:

\vspace{0.4em}
\centering
\begin{tabular}{@{}l S[table-format=1.4] S[table-format=1.4] S[table-format=1.4]
S[table-format=1.4] S[table-format=1.4] S[table-format=1.4]@{}}
\toprule
& {Random forest} & {XGBoost} & {LightGBM} & {SVM} & {QSVM angle} & {QSVM IQP} \\
\midrule
\multicolumn{7}{@{}l}{\emph{Binary}}\\
\nc{}~$=1$ & \bfseries 0.5582 & 0.5451 & 0.5093 & 0.5314 & 0.4532 & 0.4587 \\
\nc{}~$=3$ & \bfseries 0.6575 & 0.6465 & 0.6434 & 0.6306 & 0.4924 & 0.5084 \\
\nc{}~$=6$ & \bfseries 0.6835 & 0.6680 & 0.6593 & 0.6724 & 0.5138 & 0.5137 \\
\addlinespace
\multicolumn{7}{@{}l}{\emph{15-class}}\\
\nc{}~$=1$ & \bfseries 0.5562 & 0.5313 & 0.5410 & 0.5294 & 0.1688 & 0.1541 \\
\nc{}~$=3$ & 0.7194 & 0.7089 & 0.7035 & \bfseries 0.7389 & 0.3908 & 0.4797 \\
\nc{}~$=6$ & 0.7903 & 0.7659 & 0.7709 & \bfseries 0.7930 & 0.4232 & 0.5195 \\
\bottomrule
\end{tabular}

\raggedright
\take{random forest beats SVM on \emph{every} binary cell --- so the margins we
had published were measured against a sub-optimal baseline.}
\end{frame}

\begin{frame}{The margins we reported were understated}
\small
\centering
\begin{tabular}{@{}lccc@{}}
\toprule
& \nc{}~$=1$ & \nc{}~$=3$ & \nc{}~$=6$ \\
\midrule
gap vs.\ SVM (as previously reported) & $+0.0727$ & $+0.1222$ & $+0.1587$ \\
gap vs.\ \emph{best} classical & $\mathbf{+0.0995}$ & $\mathbf{+0.1491}$ & $\mathbf{+0.1697}$ \\
\midrule
understated by & 0.0268 & 0.0269 & 0.0110 \\
\bottomrule
\end{tabular}

\vspace{0.8em}
\raggedright
On the 15-class task SVM genuinely \emph{is} the strongest baseline at
\nc{}~$\in\{3,6\}$, so those two margins stand unchanged.

\vspace{0.3em}
Four of six cells move. All four move the same way.

\take{the correction goes \emph{against} the quantum pipeline everywhere. Our
earlier numbers flattered it.}
\end{frame}

\begin{frame}{Capacity scaling --- the gap widens, it does not close}
\centering
\includegraphics[width=0.96\textwidth]{\FIG{fig_capacity_scaling.pdf}}

\raggedright
\take{on EMBER binary the QSVM gains $0.055$ macro-F1 from 1 to 6 components;
the random forest gains $0.125$. More qubits are not buying the quantum model
the headroom the classical model extracts from the same components.}
\end{frame}

% ---------------------------------------------------------------- collapse
\begin{frame}{The open question we inherited}
\begin{columns}[T,onlytextwidth]
\begin{column}{0.52\textwidth}
\small
On \textbf{CIC 15-class} the QSVM sits at the random baseline
($0.056$--$0.079$ vs.\ chance $\approx 0.067$).

\vspace{0.5em}
On \textbf{EMBER 15-class} the same method reaches $0.42$--$0.52$.

\vspace{0.5em}
Already ruled out: class imbalance, qubit count.

\vspace{0.5em}
\textbf{Why?}
\end{column}
\begin{column}{0.44\textwidth}
\small
Two candidate explanations, and they are testable separately:
\begin{enumerate}\small
  \item the \textbf{features} are harder on CIC
  \item the \textbf{kernel} fails on CIC
\end{enumerate}
\vspace{0.5em}
A classical SVM extracts $0.108$--$0.165$ from \emph{the same projected
features} where the QSVM extracts $0.056$--$0.079$.
\end{column}
\end{columns}

\take{whatever CIC's features lack, the classical kernel finds something in
them that the fidelity kernel does not. So (1) alone cannot be the answer.}
\end{frame}

\begin{frame}{Feature geometry --- explains the dataset, not the kernel}
\centering
\includegraphics[width=0.90\textwidth]{\FIG{fig_geometry.pdf}}

\raggedright
\footnotesize
EMBER's families are $\approx 8\times$ more separable by Fisher ratio, up to
$50\times$ on the worst-case family. CIC's silhouette is \textbf{negative at
every} \nc{}: the average point lies closer to another family's centroid than
its own. CIC's PCA nonetheless explains \emph{more} variance ($33$--$81\%$ vs.\
$8$--$30\%$) --- its variance sits in directions that do not discriminate family.

\take{this explains why \emph{every} method does badly on CIC 15-class. It
cannot explain the QSVM-specific failure.}
\end{frame}

\begin{frame}{Kernel geometry --- the answer}
\centering
\includegraphics[width=0.90\textwidth]{\FIG{fig_kernel_diagnostics.pdf}}

\raggedright
\footnotesize
Both panels: first outer training fold, \nc{}~$=6$, IQP encoding, 800 rows. The
RBF control is built on the \textbf{identical} feature matrix --- that is what
separates ``these features are hard'' from ``this kernel is bad''.

\take{RBF produces the same spread on both datasets ($0.221$ vs.\ $0.218$). The
fidelity kernel matches it on EMBER ($0.92\times$) and \textbf{collapses to
$\mathbf{0.40\times}$ on CIC}.}
\end{frame}

\begin{frame}{What the kernel result does and does not support}
\small
\begin{block}{Supported}
The failure on CIC is a property of \textbf{the kernel}, not of the dataset.
A classical kernel on identical features does not concentrate; the quantum one
does. The Gram approaches a scaled identity --- every pair looks equally
similar --- and an SVM given that has nothing to separate on. That is what a
macro-F1 pinned to chance looks like from the inside.
\end{block}

\begin{block}{Not supported}
That tuning \texttt{bandwidth} would repair it. Bandwidth is the identified
mechanism and has never been tuned on CIC beyond its default --- but no sweep
was run. This is a hypothesis with a named lever, not a measured fix.
\end{block}

\begin{block}{One tension we state rather than resolve}
Alignment says neither kernel's Gram carries label information on CIC, yet a
classical SVM reaches $0.16$ there. Alignment is a global unweighted average
over all pairs; an SVM needs no such thing. It ranks dataset difficulty well
and predicts achievable SVM performance poorly.
\end{block}
\end{frame}

% ---------------------------------------------------------------- significance
\begin{frame}{Result 3 --- the first significance tests}
\begin{columns}[T,onlytextwidth]
\begin{column}{0.52\textwidth}
\centering
\includegraphics[width=\textwidth]{\FIG{fig_significance.pdf}}
\end{column}
\begin{column}{0.44\textwidth}
\small
Specified in Week 1. Built, unit-tested, and \textbf{never run} until now.

\vspace{0.6em}
Four weeks of ``QSVM loses to classical'' rested on differences in mean
macro-F1, with no test that any gap exceeded fold-to-fold noise.

\vspace{0.6em}
\textbf{All 84 testable pairs} favour the classical model, and all 84 reach
$p<0.05$ on the paired $t$-test.

\vspace{0.4em}
Worst $p$-value across every pair: $\mathbf{0.0241}$.
\end{column}
\end{columns}

\take{the headline claim now has evidence behind it, not an unexamined mean
difference.}
\end{frame}

\begin{frame}{Two things we will not overclaim}
\small
\begin{block}{Wilcoxon returned $0.0625$ for all 84 pairs --- that is not a null result}
With 5 paired samples, $0.0625$ is the \emph{smallest attainable} two-sided
$p$-value. The test is mathematically incapable of reaching $0.05$ at this fold
count, whatever the effect size. All 84 hitting the floor means every pair
achieved the most extreme rank configuration available.
\emph{Fix:} 6 folds would put $0.03125$ in reach --- not a different test.
\end{block}

\begin{block}{All CIC binary comparisons are untestable}
Every CIC binary run --- classical \textbf{and} quantum --- predates
parent/child run tracking and carries no recoverable sweep boundary. Raw
records show the QSVM at $0.92$--$1.00$, consistent with the earlier ``ties on
binary'' reading, but they resolve into two distinct sweeps per \nc{} with no
identifier separating them.
\emph{We did not reconstruct a grouping.} Doing so would invent a boundary the
data does not record. The limitation is symmetric across frameworks.
\end{block}
\end{frame}

% ---------------------------------------------------------------- closures
\begin{frame}{Also closed this week}
\small
\begin{block}{Subsample representativeness --- specified in Week 2, never executed}
Every quantum result uses a 1000-row subsample. Read against its null
expectation --- a faithful subsample still produces
$\approx \alpha \cdot n_{\mathrm{features}}$ rejections by chance --- CIC rejects
\textbf{nothing} of 55 features where $\approx 2.75$ were expected; EMBER rejects
\textbf{one} of 2381 where $\approx 119$ were. Class-proportion drift $< 0.001$.
\emph{Four weeks of results now rest on a measured foundation.}
\end{block}

\begin{block}{SOREL-20M --- closed, not deferred}
The labelling scheme was designed and validated over 19.7M rows and remains
delivered work; the feature store was never fetched and \textbf{no sweep was
ever run}. The two are independent --- the labelling work is not partial credit
for the sweeps. Neither disk (\SI{227}{\gibi\byte} free vs.\
\SI{71.6}{\gibi\byte}) nor memory is the blocker; it is the fixed
\SI{71.6}{\gibi\byte} transfer, because the store offers no key-level remote
access.
\end{block}
\end{frame}

% ---------------------------------------------------------------- summary
\begin{frame}{Summary}
\small
\begin{itemize}
  \item Classical models beat the QSVM on \textbf{every} testable configuration
        --- now with $p<0.05$ on all 84, not on mean differences alone.
  \item The EMBER margins were \textbf{understated}, not overstated. Every
        correction this week went against the quantum pipeline.
  \item The one configuration where the quantum pipeline was ever competitive
        --- CIC binary --- cannot be significance-tested from what was logged,
        and we say so rather than dropping it quietly.
\end{itemize}

\vspace{0.5em}
\begin{block}{The transferable result}
The fidelity kernel does not fail on CIC because CIC's features are hard --- a
classical RBF kernel on the identical features does not concentrate. It fails
because \textbf{the fidelity kernel collapses to $\mathbf{0.40\times}$ the
classical kernel's spread} on that data. That is a statement about the method,
not about the dataset.
\end{block}

\vspace{0.3em}
\textbf{Next:} tune \texttt{bandwidth} on CIC --- the named mechanism, still
untested; move to 6 folds so a non-parametric test becomes usable.
\end{frame}

\end{document}
